\section{Parsing goal and modelling stance}
The intended task is the following: given a finite corpus $\mathcal C$ of sentences and a string of
words drawn from the corpus vocabulary, we want to assign Lambek types to that string and,
when several parses are available, rank them probabilistically. In other words, the model is not
only a recogniser of grammaticality but a \emph{ranker} of candidate analyses.

The formalisation in this thesis deliberately simplifies several practical issues.
\begin{itemize}
    \item It assumes that lexical type probabilities are already available, either from manual
    annotation or from a prior learning stage.
    \item It does not handle words outside the corpus vocabulary.
    \item It treats probabilities as corpus-relative confidence values rather than as a fully learned
    generative model.
    \item It abstracts away from implementation issues such as chart parsing, estimation, and
    smoothing.
\end{itemize}
These choices are not defects so much as scope restrictions: they isolate the semantic core that
is needed for the soundness proof.

